\chapter*{Conclusions}
\addcontentsline{toc}{chapter}{Conclusions}
\label{ch:conclusions}

The Bachelor Paper presented the development of a Telegram bot information
system for public transport in Fergana city. This chapter summarizes the
work, confirms the completion of all tasks, validates the hypothesis, and
states the achievement of the aim. The chapter also describes the practical
and theoretical value of the work.

\section*{Confirmation of Tasks}
\addcontentsline{toc}{section}{Confirmation of Tasks}

Seven tasks were defined in the Introduction. Each task was completed in a
specific part of the thesis. The correspondence between the tasks and the
chapters is shown in Table~\ref{tab:tasks-conclusion}.

\begin{table}[h]
\centering
\caption{Correspondence between research tasks and thesis chapters}
\label{tab:tasks-conclusion}
\begin{tabular}{p{0.55\textwidth}l}
\hline
\textbf{Task} & \textbf{Chapter} \\
\hline
1. Review literature about information systems, Telegram bots, and
databases. & Chapter~1 \\
2. Analyze the current situation of public transport information in
Fergana. & Chapter~2 \\
3. Design the structure of the Telegram bot and the SQLite database. &
Section~3.1 \\
4. Develop the bot using Python and the aiogram library. & Section~3.1
and Section~3.2 \\
5. Add real transport data about routes, stops, and taxi services. &
Section~3.1, Section~3.2 \\
6. Test the bot with a small group of users and evaluate the results. &
Section~3.2 and Section~3.3 \\
7. Prepare proposals for future improvement of the system. & Chapter~4 \\
\hline
\end{tabular}
\end{table}

Task~1 (literature review on chat bots, Telegram Bot API, and database
design) was completed in Chapter~1, which provided the theoretical
foundation for the project. Task~2 (analysis of public transport information
in Fergana) was carried out in Chapter~2, which showed a clear gap in
available digital tools. Task~3 (database design) and Task~4 (bot
development) were accomplished and described in Sections~3.1 and~3.2.
Task~5 (data collection) was completed by importing six bus routes from
Yandex Maps into the SQLite database. Task~6 (testing and evaluation) was
carried out with 12 users from Fergana and is described in Sections~3.2
and~3.3. Task~7 (proposals for future improvement) was fulfilled in
Chapter~4, where extensions such as Uzbek language support, a button menu,
and AI integration were proposed.

\section*{Validation of the Hypothesis}
\addcontentsline{toc}{section}{Validation of the Hypothesis}

The hypothesis of the study was formulated as follows: \textit{a Telegram
bot connected to a SQLite database can give Fergana residents fast and
reliable access to transport information, and it can be more accessible
than a mobile application because Telegram is already installed on many
users' phones.}

The hypothesis was validated by the results of the user testing presented
in Chapter~3. The average response time was 0.8~seconds, which is below
the target of 2~seconds. The success rate of requests was 94\,\%. The
positive feedback rate was 92\,\% (11 of 12 test users). Test users also
confirmed that they prefer using the bot inside Telegram rather than
installing a separate mobile application. These results confirm that the
proposed solution is both fast and reliable, and that the chosen
distribution channel (Telegram) is well accepted by the target users.

\section*{Achievement of the Aim}
\addcontentsline{toc}{section}{Achievement of the Aim}

The aim of the thesis was to develop a Telegram bot information system that
helps Fergana residents and visitors find information about public transport
routes, stops, and taxi services. The aim was achieved. The developed bot,
\texttt{@FerganaTransportBot}, is deployed on the PythonAnywhere cloud
platform and is available 24/7. The bot supports six commands, contains six
bus routes with 84~stops, and was successfully tested with real users.

The three defended thesis statements presented in the Introduction were
also supported by the work:

\begin{enumerate}
    \item \textit{A new Telegram bot system for public transport
        information in Fergana city, based on Python and SQLite, was
        offered.} The bot was designed, implemented, and deployed.
    \item \textit{The traditional approach of accessing transport
        information was modified by moving it from offline sources to a
        24/7 digital channel inside Telegram.} The bot replaces the need
        to ask drivers or other people for route information.
    \item \textit{An extensible architecture that allows future
        integration with AI assistants for handling free-form questions
        was suggested.} The architecture and the proposed extension were
        described in Chapter~4.
\end{enumerate}

\section*{Practical and Theoretical Value}
\addcontentsline{toc}{section}{Practical and Theoretical Value}

\noindent\textbf{Practical value.} The Fergana Transport Bot is a real,
working system that can be used by Fergana residents and visitors. It is
free, available 24/7, and does not require the installation of a separate
mobile application. The bot can save time for users who need quick
information about bus routes, stops, or taxi services.

\noindent\textbf{Theoretical value.} The work demonstrates how a Telegram
bot can replace a simple mobile application for information delivery in
regional cities. The combination of Python, aiogram, and SQLite is shown to
be sufficient for a small-scale public transport information system. The
study also shows that a manually collected dataset can be a practical
solution when official open data is not available.

\noindent\textbf{Replicability.} The same approach can be applied to other
regional cities of Uzbekistan and to similar cities in other countries
where Telegram is widely used and official transport data is missing. The
modular database structure allows easy extension with new routes, stops,
and taxi services.

\section*{Final Statement}
\addcontentsline{toc}{section}{Final Statement}

All tasks set out in the Introduction were accomplished (Tasks~1--7
correspond to Chapters~1--4). The research hypothesis was validated, and
the aim of the thesis has been reached.